\scanpage{33}
\mspage{16}
These maps are \uline{injective}.
\deleted{\unsure{16|46}}

\uline{Proof}. It suffices to prove it for the second map, and one is reduced to the case where $X$ is \inserted{itself} an algebraic group \inserted{$G$}, extension of an abelian variety by a torus. Then the map is a homomorphism of groups
\[
(16)\qquad
\begin{array}{ccl}
G(K)&\longrightarrow&H^1(K,H_1(\overline G))\\[-.5ex]
\Vert&&\\[-.5ex]
H^0(K,G)&&
\end{array}
\]
obtained as follows. One considers for every integer $n$ the Kummer exact sequence
\[
0-{}_nG\longrightarrow G\xrightarrow{n\operatorname{id}_G}G-1
\]
whence exact cohomology sequence, which gives
\[
0-H^0(K,G)_n\longrightarrow H^1(K,{}_nG)\longrightarrow {}_nH^1(K,G)-0
\]
and passing to the limit, one finds
\[
0\longrightarrow\varprojlim_n H^0(K,G)_n\longrightarrow\varprojlim_n H^1(K,{}_nG)\longrightarrow\varprojlim_n {}_nH^1(K,G)
\]
\unsure{The|The} \unsure{composite|composite} of (16) with the \deleted{\unsure{composite}} \inserted{canonical} homomorphism
\[
H^1(K,H_{\unsure{1|n}}(\overline G))-\varprojlim_n H^1(K,{}_nG)
\]
taking into account
\[
H_{\unsure{1|n}}(\overline G)=\varprojlim_n {}_nG
\]
is none other than the homomorphism induced by
\[
H^0(K)\longrightarrow\varprojlim H^0(K,G)_n\hookrightarrow\varprojlim H^1(K,{}_nG)
\]
\margin{32}
